
\section{\BookRecordIndex{T0380}隐函数定理}\label{sec:3-5}
\zhindex{yinhanshudingle}{隐函数定理}

本节将建立\BookRecordIndex{T0380}隐函数定理的一个精细版本。粗略地说，它断言：如果 $u$ 是一条\BookRecordIndex{T0193}近似
$J$-全纯曲线，即 $\bar\partial_J(u)$ 的 $L^p$-范数充分小，并且算子 $D_u$ 满射且具有
范数一致有界的\BookRecordIndex{T0392}右逆，那么 $u$ 附近就存在真正的 \BookRecordIndex{T0003}$J$-全纯曲线。这个定理将在第 10 章的
\BookRecordIndex{T0254}粘合构造中起重要作用，因而我们会仔细说明。本节在初次阅读时可以略过。以下始终假设
$J$ 是一个固定的光滑\BookRecordIndex{T0088}殆复结构。

由命题~\ref{prop:3-1-10}，只须对某个固定的 $p>2$，在所有 $W^{1,p}$ 映射
$u:\Sigma\to M$ 所成的\BookRecordIndex{T0319}索伯列夫空间 $\mathcal B^{1,p}$ 中工作。考虑无限维巴拿赫
空间丛 $\mathcal E^p\to\mathcal B^{1,p}$，它在 $u\in\mathcal B^{1,p}$ 处的纤维是
\[
  \mathcal E_u^p
  =L^p(\Sigma,\Lambda^{0,1}\otimes_Ju^*TM),
\]
即 $\Sigma$ 上取值于拉回切丛 $u^*TM$ 的 $L^p$ 类复反线性 $1$-形式所成的空间。
非线性\BookRecordIndex{T0209}柯西--黎曼方程确定了这个丛的一个截面
$\bar\partial_J:\mathcal B^{1,p}\to\mathcal E^p$；它在一条 \BookRecordIndex{T0003}$J$-全纯曲线 $u$ 处的
\BookRecordIndex{T0310}竖直微分，就是第~\ref{sec:3-1} 节式~\eqref{eq:3-1-4} 中引入的算子
\[
  D_u:W^{1,p}(\Sigma,u^*TM)
  \longrightarrow L^p(\Sigma,\Lambda^{0,1}\otimes_Ju^*TM).
\]

现在转向算子 $D_u$ 的满射性问题。由命题~\ref{prop:3-1-11}，这等价于\BookRecordIndex{T0368}形式伴随算子
$D_u^*$ 的单射性。应用\BookRecordIndex{T0380}隐函数定理时，满射性有一个定量表述十分重要。粗略地说，这表示
某个适当\BookRecordIndex{T0392}右逆的范数不能过大。构造这种\BookRecordIndex{T0392}右逆的一种办法是取算子
\[
  Q_u=D_u^*(D_uD_u^*)^{-1}:
  L^p(\Sigma,\Lambda^{0,1}\otimes_Ju^*TM)
  \longrightarrow W^{1,p}(\Sigma,u^*TM).
\]
可以证明，形如
\[
  \lVert\eta\rVert_{W^{1,q}}
  \le c\lVert D_u^*\eta\rVert_{L^q}
\]
的不等式，会蕴含 $D_u$ 限制在 $D_u^*$ 的像上时的一个类似不等式，而这个像是
$D_u$ 的核的一个补空间。（例如见 \srcite{92, Lemma 4.5}。）构造\BookRecordIndex{T0392}右逆的另一种方法是：
通过对 $\xi$ 施加逐点条件，缩小 $D_u$ 的定义域
$W^{1,p}(\Sigma,u^*TM)$，使所得算子为双射，再以 $Q_u$ 表示这个限制算子的逆。

下一个定理是本节开头所述\BookRecordIndex{T0380}隐函数定理的精细版本。我们将证明：若 $u$ 是\BookRecordIndex{T0193}近似
$J$-全纯曲线，并且它的算子 $D_u$ 具有充分强的满射性，那么 $u$ 附近存在
\BookRecordIndex{T0003}$J$-全纯曲线，而且这些曲线可以用 $D_u$ 的核中零点的一个邻域来建模。更明确地说，对每个
充分小的截面 $\xi\in\ker D_u$，可以找到唯一一条具有形式
\[
  v_\xi=\exp_u(\xi+Q_u\eta)
\]
的 \BookRecordIndex{T0003}$J$-全纯曲线。

在\BookRecordIndex{T0254}粘合论证的应用中，允许 $\Sigma$ 上的体积形式变化至关重要。下面的注引入一类体积形式，
使我们可以在\BookRecordIndex{T0380}隐函数定理中得到一致估计。

\begin{remark}\label{rem:3-5-1}
设 $(\Sigma,j_\Sigma)$ 是带复结构的紧致二维流形。给定常数 $p>2$ 和正体积形式
$d\operatorname{vol}_\Sigma$，以 $c_p(d\operatorname{vol}_\Sigma)$ 表示索伯列夫嵌入
$W^{1,p}(\Sigma)\hookrightarrow C^0(\Sigma)$ 的范数（见定理 B.1.11）。因此
\[
  c_p(d\operatorname{vol}_\Sigma)
  :=\sup_{0\ne f\in C^\infty(\Sigma)}
  \frac{\lVert f\rVert_{L^\infty}}
       {\lVert f\rVert_{W^{1,p}}}.
\]
现在设 $E\to\Sigma$ 是黎曼向量丛，$\nabla$ 是 $E$ 上任意黎曼联络。则对每个
$W^{1,p}$ 截面 $\xi\in W^{1,p}(\Sigma,E)$，标量函数 $z\mapsto|\xi(z)|$ 属于
$W^{1,p}(\Sigma)$，并且对 $\Sigma$ 上任意度量都有
\[
  \int_\Sigma|d|\xi||^p\,d\operatorname{vol}_\Sigma
  \le
  \int_\Sigma|\nabla\xi|^p\,d\operatorname{vol}_\Sigma.
\]
要看出这一点，令函数 $z\mapsto\sqrt{|\xi(z)|^2+\varepsilon^2}$ 在
$\varepsilon\to0$ 时取极限，并使用如下事实：只要 $\xi$ 非零，就有
\[
  |d|\xi||=\frac{|\langle\nabla\xi,\xi\rangle|}{|\xi|}
  \le|\nabla\xi|.
\]
由此
\[
  \sup_\Sigma|\xi|
  \le c_p(d\operatorname{vol}_\Sigma)
  \left(\int_\Sigma(|\xi|^p+|\nabla\xi|^p)\,
  d\operatorname{vol}_\Sigma\right)^{1/p}
\]
对每个黎曼向量丛 $E\to\Sigma$ 的每个 $W^{1,p}$ 截面 $\xi$ 以及每个黎曼联络
$\nabla$ 都成立。
\end{remark}

在下一个定理中，固定带复结构的紧致二维流形 $(\Sigma,j_\Sigma)$，并考虑满足
$c_p(d\operatorname{vol}_\Sigma)$ 具有一致上界的所有体积形式
$d\operatorname{vol}_\Sigma$。体积形式 $d\operatorname{vol}_\Sigma$ 与复结构
$j_\Sigma$ 一起确定 $\Sigma$ 上的度量，以下所有相关范数都关于这个度量理解。丛
$u^*TM\to\Sigma$ 截面的索伯列夫范数，则关于由 $\omega$ 和 $J$ 确定的度量的
\BookRecordIndex{T0034}Levi-Civita 联络 $\nabla$ 来理解。

\begin{theorem}\label{thm:3-5-2}
固定常数 $p>2$ 和闭黎曼曲面 $(\Sigma,j_\Sigma)$。则对每个常数 $c_0>0$，都存在常数
$\delta>0$，使下述结论对 $\Sigma$ 上每个满足
$c_p(d\operatorname{vol}_\Sigma)\le c_0$ 的体积形式
$d\operatorname{vol}_\Sigma$ 成立。假设
$u\in W^{1,p}(\Sigma,M)$ 和 $\xi_0\in W^{1,p}(\Sigma,u^*TM)$ 满足
\[
  \lVert du\rVert_{L^p}\le c_0,\qquad
  \lVert\xi_0\rVert_{W^{1,p}}<\frac\delta8,\qquad
  \bigl\lVert\bar\partial_J(\exp_u(\xi_0))\bigr\rVert_{L^p}
  <\frac\delta{4c_0}.
\]
再假设
\[
  Q_u:L^p(\Sigma,\Lambda^{0,1}\otimes_Ju^*TM)
  \longrightarrow W^{1,p}(\Sigma,u^*TM)
\]
是 $D_u$ 的\BookRecordIndex{T0392}右逆，并且
\[
  D_uQ_u=\mathbf1,\qquad \lVert Q_u\rVert\le c_0.
\]
则存在唯一的截面 $\xi=Q_u\eta\in W^{1,p}(\Sigma,u^*TM)$，使得
\[
  \bar\partial_J(\exp_u(\xi_0+\xi))=0,\qquad
  \lVert\xi_0+\xi\rVert_{W^{1,p}}<\delta.
\]
此外，
\[
  \lVert\xi\rVert_{W^{1,p}}
  \le2c_0\bigl\lVert\bar\partial_J(\exp_u(\xi_0))\bigr\rVert_{L^p}.
\]
\end{theorem}

证明是\BookRecordIndex{T0380}隐函数定理的一个应用。考虑空间
\[
  \mathcal X_u:=W^{1,p}(\Sigma,u^*TM),
  \qquad
  \mathcal Y_u:=L^p(\Sigma,\Lambda^{0,1}\otimes_Ju^*TM),
\]
以及式~\eqref{eq:3-1-3} 定义的映射
$\mathcal F_u:\mathcal X_u\to\mathcal Y_u$，即
\[
  \mathcal F_u(\xi)
  :=\Phi_u(\xi)^{-1}\bigl(\bar\partial_J(\exp_u(\xi))\bigr),
\]
其中 $\xi\in\mathcal X_u$，而
$\Phi_u(\xi):u^*TM\to\exp_u(\xi)^*TM$ 由沿测地线
$s\mapsto\exp_{u(z)}(s\xi(z))$ 关于\BookRecordIndex{T0052}埃尔米特联络 $\widetilde\nabla$ 的\BookRecordIndex{T0263}平行移动定义。由
命题~\ref{prop:3-1-1}，$\mathcal F_u$ 在零点的微分为
\[
  d\mathcal F_u(0)=D_u.
\]
下一个结果可以解释为 $\mathcal F_u$ 的一个\BookRecordIndex{T0115}二次估计（注~\ref{rem:3-5-5}）。我们把它
表述为映射
$\mathcal F_u:\mathcal X_u\to\mathcal Y_u$
\sourcefootnote{原书此处把映射误写作
$\mathcal F_u(\xi):\mathcal X_u\to\mathcal Y_u$；但
$\mathcal F_u(\xi)$ 是 $\mathcal Y_u$ 中的取值，而紧邻上文定义的映射是
$\mathcal F_u$，故据源靶类型改正。}的微分
$d\mathcal F_u(\xi):\mathcal X_u\to\mathcal Y_u$ 的估计。注意，$\Sigma$ 上的体积形式
不参与映射 $\mathcal F_u$ 的定义。

\begin{proposition}\label{prop:3-5-3}
固定常数 $p>2$ 和闭黎曼曲面 $(\Sigma,j_\Sigma)$。则对每个常数 $c_0>0$，都存在常数
$c>0$，使下述结论对 $\Sigma$ 上每个满足
$c_p(d\operatorname{vol}_\Sigma)\le c_0$ 的体积形式
$d\operatorname{vol}_\Sigma$ 成立。若
$u\in W^{1,p}(\Sigma,M)$ 和 $\xi\in W^{1,p}(\Sigma,u^*TM)$ 满足
\begin{equation}\label{eq:3-5-1}
  \lVert du\rVert_{L^p}\le c_0,\qquad
  \lVert\xi\rVert_{L^\infty}\le c_0,
\end{equation}
则
\begin{equation}\label{eq:3-5-2}
  \lVert\!\lVert\!\lVert
  d\mathcal F_u(\xi)-D_u
  \rVert\!\rVert\!\rVert
  \le c\lVert\xi\rVert_{W^{1,p}}.
\end{equation}
这里三重竖线表示 $\mathcal L(\mathcal X_u,\mathcal Y_u)$ 上的算子范数。
\end{proposition}

\begin{proof}
证明通过直接计算得到。给定 $x\in M$ 和 $\xi\in T_xM$，定义（双）线性映射
\[
  E_x(\xi):T_xM\longrightarrow T_{\exp_x(\xi)}M,
  \qquad
  \Psi_x(\xi):T_xM\times T_xM\longrightarrow T_{\exp_x(\xi)}M
\]
如下。以 $\Phi_x$ 表示上面\BookRecordIndex{T0263}平行移动的逐点版本
\sourcefootnote{原书在 $\Psi_x$ 的点态定义及紧随其后的说明中把 $\Phi_x$ 两次印作
$\Phi_u$；但此处给定的是 $x\in M$，相关测地线以 $x$ 为基点，且源靶均以 $x$ 标记，
故据类型以及随后从逐点 $\Psi_x$ 得到 $\Psi_u$ 的记号唯一改为 $\Phi_x$。}，令
\[
  E_x(\xi)\xi'
  :=\left.\frac d{dt}\right|_{t=0}\exp_x(\xi+t\xi'),
  \qquad
  \Psi_x(\xi;\xi',\eta)
  :=\left.\widetilde\nabla_t
  \bigl(\Phi_x(\xi+t\xi')\eta\bigr)\right|_{t=0}.
\]
注意，$\Psi_x(0;\xi',\eta)=0$，因为当 $\xi=0$ 时，这里沿定义
$\Phi_x(t\xi')$ 时所用的同一条测地线取协变导数。

现在对恒等式
\[
  \Phi_u(\xi+t\xi')\mathcal F_u(\xi+t\xi')
  =\bar\partial_J(\exp_u(\xi+t\xi'))
\]
在 $t=0$ 处关于 $\widetilde\nabla$ 作协变微分，得到
\[
  \Phi_u(\xi)d\mathcal F_u(\xi)\xi'
  +\Psi_u(\xi;\xi',\mathcal F_u(\xi))
  =D_{\exp_u(\xi)}\bigl(E_u(\xi)\xi'\bigr).
\]
把这个恒等式写成
\begin{equation}\label{eq:3-5-3}
  d\mathcal F_u(\xi)\xi'
  =\Phi_u(\xi)^{-1}D_{\exp_u(\xi)}\bigl(E_u(\xi)\xi'\bigr)
  -\Phi_u(\xi)^{-1}\Psi_u(\xi;\xi',\mathcal F_u(\xi)).
\end{equation}
选取常数 $c_1>0$，使得对每个 $x\in M$、每个满足 $|\xi|\le c_0$ 的
$\xi\in T_xM$ 以及每个 $\eta\in T_xM$，都有
\[
  |E_x(\xi)|\le c_1,\qquad
  |\Psi_x(\xi;\xi',\eta)|\le c_1|\xi||\xi'||\eta|.
\]
在适当之处，$|\cdot|$ 表示
$\mathcal L(T_xM,T_{\exp_x(\xi)}M)$ 上的算子范数；第二个不等式成立是因为
$\Psi_x(0;\xi',\eta)=0$。于是式~\eqref{eq:3-5-3} 的第二项满足逐点估计
\[
  \bigl|\Phi_u(\xi)^{-1}
  \Psi_u(\xi;\xi',\mathcal F_u(\xi))\bigr|
  \le c_1|d\exp_u(\xi)|\,|\xi|\,|\xi'|,
\]
其中 $d\exp_u(\xi)$ 是映射
$\exp_u(\xi):\Sigma\to M$
\sourcefootnote{原书把这个映射的定义域印作 $S^2$；但本节从始至终固定的是任意闭黎曼曲面
$\Sigma$，且 $\exp_u(\xi)$ 是逐点指数映射所得的 $\Sigma\to M$ 映射，故改为
$\Sigma$。}的微分。存在只依赖于 $c_0$ 和 $M$ 上度量的常数 $c_2$，使得只要
$u,\xi$ 满足式~\eqref{eq:3-5-1}，就有
\[
  |d\exp_u(\xi)|\le c_2(|du|+|\nabla\xi|).
\]
因此
\[
  \bigl\lVert\Phi_u(\xi)^{-1}
  \Psi_u(\xi;\xi',\mathcal F_u(\xi))\bigr\rVert_{L^p}
  \le c_1c_2
  \bigl(\lVert du\rVert_{L^p}+\lVert\nabla\xi\rVert_{L^p}\bigr)
  \lVert\xi\rVert_{L^\infty}\lVert\xi'\rVert_{L^\infty}.
\]
由于
$\lVert du\rVert_{L^p}\le c_0$、
$\lVert\xi\rVert_{L^\infty}\le c_0$，并且
$\lVert\xi'\rVert_{L^\infty}\le
c_0\lVert\xi'\rVert_{W^{1,p}}$
\sourcefootnote{原书在这个索伯列夫嵌入估计中漏印了 $\xi'$ 的撇号，因而把前一项
$\xi$ 的估计重复了一次。上一行含有因子 $\lVert\xi'\rVert_{L^\infty}$，而紧随其后的
式~\eqref{eq:3-5-4} 要把它估计为 $\lVert\xi'\rVert_{W^{1,p}}$；常数
$c_3=(c_0+1)c_0^2c_1c_2$ 也印证了这个唯一补正。}，可得
\begin{equation}\label{eq:3-5-4}
  \bigl\lVert\Phi_u(\xi)^{-1}
  \Psi_u(\xi;\xi',\mathcal F_u(\xi))\bigr\rVert_{L^p}
  \le c_3\lVert\xi\rVert_{W^{1,p}}
  \lVert\xi'\rVert_{W^{1,p}},
\end{equation}
其中 $c_3:=(c_0+1)c_0^2c_1c_2$。

下面说明式~\eqref{eq:3-5-3} 的第一项可以作如下估计：存在常数 $c_4>0$，它只依赖
$c_0$（以及 $J,\omega,j_\Sigma$），使得只要 $u,\xi$ 满足式~\eqref{eq:3-5-1}，就有
\begin{equation}\label{eq:3-5-5}
  \left\lVert
  \Phi_u(\xi)^{-1}D_{\exp_u(\xi)}\bigl(E_u(\xi)\xi'\bigr)-D_u\xi'
  \right\rVert_{L^p}
  \le c_4\lVert\xi\rVert_{W^{1,p}}
  \lVert\xi'\rVert_{W^{1,p}}.
\end{equation}
估计~\eqref{eq:3-5-4} 与~\eqref{eq:3-5-5} 立即蕴含所需结论~\eqref{eq:3-5-2}。

所以只须证明式~\eqref{eq:3-5-5}。记
$u_\xi:=\exp_u(\xi)$，并考察恒等式
\begin{align*}
  D_{u_\xi}\bigl(E_u(\xi)\xi'\bigr)-\Phi_u(\xi)D_u\xi'
  ={}&\bigl(\nabla(E_u(\xi)\xi')-
  \Phi_u(\xi)\nabla\xi'\bigr)^{0,1}\\
  &-\frac12J(u_\xi)
  \Bigl(
    (\nabla_{E_u(\xi)\xi'}J)(u_\xi)du_\xi
    -\Phi_u(\xi)(\nabla_{\xi'}J)(u)du
  \Bigr)^{0,1}.
\end{align*}
分别考察右端的两个加项。由于 $J$ 是 $TM$ 的等距映射，第二个加项中的各项逐点估计为
\begin{align*}
  &\bigl|
  (\nabla_{E_u(\xi)\xi'}J)(u_\xi)du_\xi
  -\Phi_u(\xi)(\nabla_{\xi'}J)(u)du
  \bigr|\\
  &\quad\le
  \bigl|
  (\nabla_{E_u(\xi)\xi'}J)(u_\xi)du_\xi
  -(\nabla_{E_u(\xi)\xi'}J)(u_\xi)\Phi_u(\xi)du
  \bigr|\\
  &\qquad+
  \bigl|
  (\nabla_{E_u(\xi)\xi'}J)(u_\xi)\Phi_u(\xi)du
  -\Phi_u(\xi)(\nabla_{\xi'}J)(u)du
  \bigr|\\
  &\quad\le
  c_1\lVert\nabla J\rVert_{L^\infty}|\xi'|
  |du_\xi-\Phi_u(\xi)du|
  +c_5|du|\,|\xi|\,|\xi'|\\
  &\quad\le c_6\bigl(|du|\,|\xi|+|\nabla\xi|\bigr)|\xi'|.
\end{align*}
这里常数 $c_6$ 只依赖于 $c_0$ 和 $M$ 上的度量。第一个加项的逐点估计具有形式
\begin{align*}
  \bigl|\nabla(E_u(\xi)\xi')-\Phi_u(\xi)\nabla\xi'\bigr|
  &\le
  \bigl|(E_u(\xi)-\Phi_u(\xi))\nabla\xi'\bigr|
  +\bigl|\nabla(E_u(\xi)\xi')-E_u(\xi)\nabla\xi'\bigr|\\
  &\le c_7\bigl(
  |\xi|\,|\nabla\xi'|+|du|\,|\xi|\,|\xi'|
  +|\nabla\xi|\,|\xi'|\bigr).
\end{align*}
合并这些估计，得到逐点不等式
\[
  \bigl|D_{u_\xi}(E_u(\xi)\xi')-\Phi_u(\xi)D_u\xi'\bigr|
  \le(c_6+c_7)\bigl(
  |\xi|\,|\nabla\xi'|+|du|\,|\xi|\,|\xi'|
  +|\nabla\xi|\,|\xi'|\bigr).
\]
现在取 $p$ 次方并积分，分别对因子 $du,\nabla\xi,\nabla\xi'$ 使用 $L^p$-范数，
对其他因子使用 $L^\infty$-范数。由于
\[
  \lVert\xi\rVert_{L^\infty}
  \le c_0\lVert\xi\rVert_{W^{1,p}},
\]
这就证明了式~\eqref{eq:3-5-5} 和命题~\ref{prop:3-5-3}。
\end{proof}

\begin{proof}[定理~\ref{thm:3-5-2} 的证明]
根据假设，算子 $D_u=d\mathcal F_u(0)$ 满射，并且有\BookRecordIndex{T0392}右逆
$Q_u:\mathcal Y_u\to\mathcal X_u$，满足
\[
  \lVert Q_u\rVert\le c_0.
\]
令 $c$ 为命题~\ref{prop:3-5-3} 中的常数，并选取 $\delta\in(0,1)$，使得
\[
  c\delta\le\frac1{2c_0}.
\]
于是由式~\eqref{eq:3-5-2}，
\[
  \lVert\xi\rVert_{W^{1,p}}<\delta
  \quad\Longrightarrow\quad
  \lVert\!\lVert\!\lVert
  d\mathcal F_u(\xi)-D_u
  \rVert\!\rVert\!\rVert
  \le\frac1{2c_0}.
\]
因此命题 A.3.4 的假设对
$X=\mathcal X_u$、$Y=\mathcal Y_u$、$f=\mathcal F_u$、$x_0=0$、$c=c_0$ 以及这个常数
$\delta$ 成立。所以，如果 $\xi_0\in\mathcal X_u$ 满足
\[
  \lVert\xi_0\rVert_{W^{1,p}}<\frac\delta8,
  \qquad
  \lVert\mathcal F_u(\xi_0)\rVert_{L^p}
  =\bigl\lVert\bar\partial_J(\exp_u(\xi_0))\bigr\rVert_{L^p}
  <\frac\delta{4c_0},
\]
那么由命题 A.3.4，存在唯一的 $\xi\in\operatorname{im}Q_u$，使得
\[
  \mathcal F_u(\xi_0+\xi)=0,
  \qquad
  \lVert\xi_0+\xi\rVert_{W^{1,p}}<\delta.
\]
此外，这个向量 $\xi$ 满足
\[
  \lVert\xi\rVert_{W^{1,p}}
  \le2c_0\lVert\mathcal F_u(\xi_0)\rVert_{L^p}
  =2c_0\bigl\lVert\bar\partial_J(\exp_u(\xi_0))\bigr\rVert_{L^p}.
\]
这就证明了定理~\ref{thm:3-5-2}。
\end{proof}

\begin{remark}\label{rem:3-5-4}
定理~\ref{thm:3-5-2} 是对固定\BookRecordIndex{T0088}殆复结构 $J\in\mathcal J_\tau(M,\omega)$ 陈述的。
第 10 章需要该定理的一个版本，其中\BookRecordIndex{T0088}殆复结构依赖于基点 $z\in\Sigma$。（见式 (6.7.4)
以及第 6.7 节的注 6.7.3。）固定一个定义在紧致流形 $\Lambda$ 上的光滑映射
\[
  \Lambda\longrightarrow\mathcal J_\tau(M,\omega),
  \qquad \lambda\longmapsto J_\lambda.
\]
我们希望修改定理~\ref{thm:3-5-2}，使其适用于具有形式
$J_z:=J_{\lambda(z)}$ 的光滑\BookRecordIndex{T0088}殆复结构族，其中 $\lambda:\Sigma\to\Lambda$ 是光滑映射。
关键在于，映射 $\lambda:\Sigma\to\Lambda$ 的导数并不进入命题~\ref{prop:3-5-3}
的\BookRecordIndex{T0115}二次估计。因此，定理~\ref{thm:3-5-2} 的断言继续成立：对每个光滑映射
$\lambda:\Sigma\to\Lambda$ 以及每条具有相应\BookRecordIndex{T0392}右逆 $Q_u$ 的 $J_\lambda$-全纯曲线
$u:\Sigma\to M$，常数 $c_0$ 和 $\delta$ 可以一致选取。
\end{remark}

\begin{remark}\label{rem:3-5-5}
命题~\ref{prop:3-5-3} 的证明表明，只要 $u,\xi$ 满足式~\eqref{eq:3-5-1}，就有估计
\[
  \lVert d\mathcal F_u(\xi)\xi'-D_u\xi'\rVert_{L^p}
  \le c\bigl(
  \lVert\xi\rVert_{L^\infty}\lVert\xi'\rVert_{W^{1,p}}
  +\lVert\xi\rVert_{W^{1,p}}\lVert\xi'\rVert_{L^\infty}
  \bigr).
\]
稍微复杂一些的论证表明，只要 $u,\xi_0,\xi_1$ 满足式~\eqref{eq:3-5-1}，就有
\[
  \lVert d\mathcal F_u(\xi_1)\xi-d\mathcal F_u(\xi_0)\xi\rVert_{L^p}
  \le c\bigl(
  \lVert\xi_1-\xi_0\rVert_{L^\infty}\lVert\xi\rVert_{W^{1,p}}
  +\lVert\xi_1-\xi_0\rVert_{W^{1,p}}\lVert\xi\rVert_{L^\infty}
  \bigr).
\]
积分后得到\BookRecordIndex{T0115}二次估计
\[
  \lVert\mathcal F_u(\xi_0+\xi)-\mathcal F_u(\xi_0)
  -d\mathcal F_u(\xi_0)\xi\rVert_{L^p}
  \le c\lVert\xi\rVert_{L^\infty}\lVert\xi\rVert_{W^{1,p}},
\]
其中 $u,\xi_0,\xi$ 满足式~\eqref{eq:3-5-1}。所有这些估计对于 $\Sigma$ 上满足
$c_p(d\operatorname{vol}_\Sigma)\le c_0$ 的所有体积形式都一致成立。
\end{remark}

使用注~\ref{rem:3-5-5} 中的估计，可以把定理~\ref{thm:3-5-2} 的唯一性陈述精细化如下。
这个精细版本在本书中不会使用。

\begin{corollary}\label{cor:3-5-6}
设 $(M,\omega,J)$、$(\Sigma,j_\Sigma)$ 与 $p>2$ 如定理~\ref{thm:3-5-2} 中所述。
则对每个 $c_0>0$，都存在 $\delta>0$，使下述结论对每个满足
$c_p(d\operatorname{vol}_\Sigma)\le c_0$ 的体积形式
$d\operatorname{vol}_\Sigma$ 成立。若
$u\in W^{1,p}(\Sigma,M)$、
$\xi_0,\xi_1\in W^{1,p}(\Sigma,u^*TM)$，而 $D_u$ 有\BookRecordIndex{T0392}右逆 $Q_u$，并且
\[
  \bar\partial_J(\exp_u(\xi_0))=0,\qquad
  \bar\partial_J(\exp_u(\xi_1))=0,\qquad
  \xi_1-\xi_0\in\operatorname{im}Q_u,
\]
\[
  \lVert du\rVert_{L^p}\le c_0,\qquad
  \lVert Q_u\rVert\le c_0,\qquad
  \lVert\xi_0\rVert_{W^{1,p}}\le\delta,\qquad
  \lVert\xi_1\rVert_{L^\infty}\le\delta,
\]
则 $\xi_0=\xi_1$。
\end{corollary}

\begin{proof}
记 $\xi:=\xi_1-\xi_0$，于是
$\xi=Q_uD_u\xi$，且
$\mathcal F_u(\xi_0)=\mathcal F_u(\xi_0+\xi)=0$。此外，
\[
  \lVert\xi\rVert_{L^\infty}
  \le\lVert\xi_0\rVert_{L^\infty}
  +\lVert\xi_1\rVert_{L^\infty}
  \le(c_0+1)\delta.
\]
选取 $c>0$，使命题~\ref{prop:3-5-3} 与注~\ref{rem:3-5-5} 中的估计成立。则
\begin{align*}
  \lVert\xi\rVert_{W^{1,p}}
  &\le c_0\lVert D_u\xi\rVert_{L^p}\\
  &\le c_0\lVert
    \mathcal F_u(\xi_0+\xi)-\mathcal F_u(\xi_0)
    -d\mathcal F_u(\xi_0)\xi
    \rVert_{L^p}
    +c_0\lVert(d\mathcal F_u(\xi_0)-D_u)\xi\rVert_{L^p}\\
  &\le c_0c\bigl(
    \lVert\xi\rVert_{L^\infty}
    +\lVert\xi_0\rVert_{W^{1,p}}
    \bigr)\lVert\xi\rVert_{W^{1,p}}\\
  &\le(c_0+2)c_0c\delta\lVert\xi\rVert_{W^{1,p}}.
\end{align*}
若 $(c_0+2)c_0c\delta<1$，就有 $\xi=0$。这证明了推论~\ref{cor:3-5-6}。
\end{proof}
